\section{Proposed Method: \diffgnnorder}
\label{sec:method}
This paper treats \diffgnnorder as the proposed method. In the codebase, it corresponds to the \texttt{diff\_gnn\_order} implementation. The placement-only \diffgnn model is a reduced variant that removes the ordering heads and is considered only in ablation studies.

\subsection{Input Feature Preparation}
\paragraph{Node features.}
For each node $i \in \gV$, the raw node feature vector is constructed from execution, area, and graph-structural information:
\begin{equation}
\vf_i^{\mathrm{raw}}=
\left[
t_i^{sw},\;
t_i^{hw},\;
A_i,\;
d_i,\;
\frac{t_i^{hw}}{t_i^{sw}+\epsilon},\;
t_i^{sw}-t_i^{hw},\;
\frac{t_i^{sw}-t_i^{hw}}{\max(A_i,\epsilon)}
\right].
\end{equation}
The first five terms are the default execution and resource descriptors; the last two correspond to the hardware-gain features inspired by \gcps \cite{zheng2021hardware}. In the repository, the strongest order-aware configuration uses the hybrid \texttt{default\_plus\_paper} profile, which concatenates both descriptor sets after normalization.

\paragraph{Edge features.}
For each edge $(j,i)\in\gE$, the raw edge feature vector is
\begin{equation}
\ve_{ji}^{\mathrm{raw}}=
\left[
C_{ji},\;
|g_j-g_i|,\;
|ga_j-ga_i|
\right],
\end{equation}
where $g_i=t_i^{sw}-t_i^{hw}$ and $ga_i=g_i/\max(A_i,\epsilon)$. These features are max-normalized and used by the optional learned edge-weight module.

\paragraph{Edge-weight mechanism.}
The code supports both static and learned edge weighting. The static component can be derived from communication magnitude, cosine similarity, or hybrid variants, while the learned component is either a per-edge parameter or an \mlp applied to edge features. In the common \diffgnnorder configuration, the base edge mode is \texttt{paper2\_cosine} and the learned edge-weight mechanism is \texttt{mlp}, with edge-weight scales clipped to $[0.5,1.5]$.

\subsection{Order-Aware Graph Encoder}
Let $h_i^{(0)}$ be the normalized input feature for node $i$. A shared \gcn-based encoder with $L$ layers computes node embeddings
\begin{equation}
h_i^{(\ell+1)}=
\mathrm{GNN}^{(\ell)}\!\left(
h_i^{(\ell)},
\{h_j^{(\ell)}:(j,i)\in\gE\},
e_{ji}
\right),
\qquad \ell=0,\dots,L-1.
\end{equation}
The default order-aware model in the repository uses a three-layer encoder with hidden dimension 256 and dropout 0.2 on the large SqueezeNet benchmark; the smaller 11-node configuration uses the same structure with hidden dimension 128 for faster sweeps.

The final embedding is passed to three heads:
\begin{equation}
\tilde{p}_i = \sigma\!\left(\mathrm{Head}_{\mathrm{place}}(h_i^{(L)})\right),
\end{equation}
\begin{equation}
\phi_i^{sw} = \mathrm{Head}_{\mathrm{sw}}(h_i^{(L)}),
\qquad
\phi_i^{hw} = \mathrm{Head}_{\mathrm{hw}}(h_i^{(L)}).
\end{equation}
The placement head predicts the relaxed hardware assignment. The software-order head predicts the software execution priority. The hardware-order head is implemented in the code, but the stable configuration used in this paper keeps \texttt{use\_hw\_ordering=False}; consequently, software ordering is the active learned ordering signal.

\subsection{Soft Ordering and Differentiable Scheduler}
\paragraph{Soft software permutation.}
The software-priority vector is converted into a soft permutation through Sinkhorn normalization:
\begin{equation}
\mP^{sw}=\mathrm{Sinkhorn}\!\left(\frac{\phi^{sw}}{\tau_o}\right),
\label{eq:sinkhorn_sw}
\end{equation}
where $\tau_o$ is the ordering temperature. We use 12 Sinkhorn iterations in the repository's fast and default order-aware settings, and Gumbel noise is disabled in the stable benchmark configuration.

\paragraph{Pairwise-before relation.}
The pairwise ordering matrix is approximated from expected ranks:
\begin{equation}
\mB^{sw}_{ji}
=
\sigma\!\left(\frac{\tilde{r}_i-\tilde{r}_j}{\tau_r}\right),
\qquad
\tilde{r}_i=\sum_{a=0}^{|\gV|-1} a \, \mP^{sw}_{ai},
\end{equation}
where $\tau_r$ is the pairwise-temperature parameter. The code uses \texttt{pairwise\_mode=rank\_sigmoid} with $\tau_r=0.35$.

\paragraph{Placement-gated software serialization.}
Ordering matters only when both tasks are assigned to software. Therefore, the effective software-order matrix is
\begin{equation}
\mR^{sw}_{ji}=
\mB^{sw}_{ji}(1-\tilde{p}_j)(1-\tilde{p}_i).
\label{eq:effective_sw_order}
\end{equation}
The relaxed DAG-readiness term follows Equation~\eqref{eq:soft_rdag}. The software-readiness term uses the learned order:
\begin{equation}
s_i^{sw}
=
\SMax_{\beta}
\Big(
\{\, \tilde{f}_j + \alpha \log(\mR^{sw}_{ji}+\epsilon) \mid j \neq i\,\}
\Big),
\label{eq:sw_readiness_order}
\end{equation}
where $\alpha$ is the resource-order scaling factor.

The relaxed start and finish times are then computed as in Equations~\eqref{eq:soft_start} and \eqref{eq:soft_finish}. This yields the order-aware relaxed makespan $\widetilde{M}(\tilde{\vp},\tilde{\vpi}^{sw})$.

\subsection{Training Objective}
The training loss combines the relaxed makespan surrogate with feasibility and regularization terms:
\begin{equation}
\mathcal{L}_{\mathrm{ord}}
=
\widetilde{M}(\tilde{\vp},\tilde{\vpi}^{sw})
+ \lambda_a \mathcal{L}_{\mathrm{area}}(\tilde{\vp})
+ \lambda_c \widetilde{C}(\tilde{\vp})
+ \lambda_{\mathrm{perm}} \Omega(\mP^{sw}),
\label{eq:ordered_loss}
\end{equation}
where $\widetilde{C}(\tilde{\vp})$ is the relaxed partition-cost term and $\Omega(\mP^{sw})$ contains optional doubly stochastic and entropy regularizers. In the benchmark configuration used here, the partition-cost coefficient is set to 0.01 on the small graph and 0.0833333333 in the default SqueezeNet configuration, while permutation regularizers are disabled.

\subsection{Discrete Decode, Repair, and Reporting}
At inference time, the relaxed placement is discretized, repaired, and post-processed:
\begin{enumerate}
\item threshold the relaxed placement scores and build one or more decode candidates,
\item repair any area-constraint violation by moving selected tasks back to software,
\item optionally refill unused hardware area according to the decode score,
\item apply hybrid LSSP-based local search.
\end{enumerate}
The hybrid post-process used in our experiments employs LSSP evaluation, area fill, node swapping, two DLS refinement steps, and a maximum of 120 local-search iterations.

We compute three schedule metrics for the repaired partition:
\begin{equation}
M_{\mathrm{DAG}},
\qquad
M_{\mathrm{static}},
\qquad
M_{\mathrm{learned}}.
\end{equation}
$M_{\mathrm{DAG}}$ is the DAG makespan without the learned software-priority signal. $M_{\mathrm{static}}$ is the LSSP makespan under static priorities computed from the repaired partition. $M_{\mathrm{learned}}$ is the LSSP makespan under the learned software-priority scores and is only defined for \diffgnnorder.

The reported value is
\begin{equation}
M_{\mathrm{report}}=
\begin{cases}
M_{\mathrm{static}}, & \text{for random, greedy, \gcps, PSO, DBPSO, CLPSO, CCPSO, ESA, SHADE, JADE, GL25, and \diffgnn,}\\
\min(M_{\mathrm{static}}, M_{\mathrm{learned}}), & \text{for \diffgnnorder.}
\end{cases}
\label{eq:report_metric}
\end{equation}
We still list $M_{\mathrm{DAG}}$ in the results tables because it provides a schedule baseline shared by all methods.

\subsection{Relation to the Placement-Only Ablation}
The placement-only \diffgnn ablation uses the same graph encoder and relaxed placement mechanism but removes the ordering heads and the order-aware scheduling terms. Its reported result is therefore based on $M_{\mathrm{static}}$ only. This ablation isolates the value of learned ordering and is the primary within-family comparison in the experimental section.
